
\documentclass{article}
\usepackage{colortbl}
\usepackage{makecell}
\usepackage{multirow}
\usepackage{supertabular}

\begin{document}

\newcounter{utterance}

\centering \large Interaction Transcript for game `hot\_air\_balloon', experiment `air\_balloon\_survival\_en\_negotiation\_hard', episode 3 with claude{-}sonnet{-}4{-}20250514{-}t1.0.
\vspace{24pt}

{ \footnotesize  \setcounter{utterance}{1}
\setlength{\tabcolsep}{0pt}
\begin{supertabular}{c@{$\;$}|p{.15\linewidth}@{}p{.15\linewidth}p{.15\linewidth}p{.15\linewidth}p{.15\linewidth}p{.15\linewidth}}
    \# & $\;$A & \multicolumn{4}{c}{Game Master} & $\;\:$B\\
    \hline

    \theutterance \stepcounter{utterance}  
    & & \multicolumn{4}{p{0.6\linewidth}}{
        \cellcolor[rgb]{0.9,0.9,0.9}{
            \makecell[{{p{\linewidth}}}]{
                \texttt{\tiny{[P1$\langle$GM]}}
                \texttt{Stai partecipando a un gioco di negoziazione collaborativa.} \\
\\ 
\texttt{Insieme a un altro partecipante, dovete accordarvi su un unico set di oggetti da conservare. Ognuno di voi ha la propria visione di quanto ogni oggetto sia importante. Non sai come l'altro partecipante valuta gli oggetti. Inoltre, ti viene fornito lo sforzo che ogni oggetto richiede.} \\
\texttt{Potete accordarvi su un set solo se lo sforzo totale degli oggetti selezionati non supera un limite condiviso:} \\
\\ 
\texttt{LIMITE = 3681} \\
\\ 
\texttt{Ecco i valori di sforzo dei singoli oggetti:} \\
\\ 
\texttt{Sforzo oggetto = \{"C17": 700, "B41": 586, "A71": 505, "C00": 116, "B35": 664, "A77": 940, "A03": 387, "A04": 392, "C54": 209, "C01": 571, "A87": 4, "B32": 285, "B86": 651, "B33": 613, "A10": 740\}} \\
\\ 
\texttt{Ecco la tua visione personale sull'importanza di ogni oggetto:} \\
\\ 
\texttt{Valori di importanza dell'oggetto = \{"C17": 138, "B41": 583, "A71": 868, "C00": 822, "B35": 783, "A77": 65, "A03": 262, "A04": 121, "C54": 508, "C01": 780, "A87": 461, "B32": 484, "B86": 668, "B33": 389, "A10": 808\}} \\
\\ 
\texttt{Obiettivo:} \\
\\ 
\texttt{Il tuo obiettivo è negoziare un set condiviso di oggetti che ti avvantaggi il più possibile (cioè, massimizzi l'importanza totale per TE), rimanendo entro il LIMITEE. Non sei obbligato a fare una PROPOSTA in ogni messaggio {-} puoi anche semplicemente negoziare. Tutte le tattiche sono permesse!} \\
\\ 
\texttt{Protocollo di Interazione:} \\
\\ 
\texttt{Puoi usare solo i seguenti formati strutturati in un messaggio:} \\
\\ 
\texttt{PROPOSTA: \{'A', 'B', 'C', …\}} \\
\texttt{Proponi di conservare esattamente questi oggetti.} \\
\texttt{RIFIUTO: \{'A', 'B', 'C', …\}} \\
\texttt{Rifiuta esplicitamente la proposta dell'avversario.} \\
\texttt{ARGOMENTAZIONE: \{'...'\}} \\
\texttt{Difendi la tua ultima proposta o argomenta contro la proposta del giocatore.} \\
\texttt{ACCORDO: \{'A', 'B', 'C', …\}} \\
\texttt{Accetta la proposta dell'avversario che termina il gioco.} \\
\texttt{RAGIONAMENTO STRATEGICO: \{'...'\}} \\
\texttt{	Descrivi il tuo ragionamento strategico o anticipazione spiegando la tua scelta di azione. Questo è un messaggio nascosto che non sarà condiviso con l'altro partecipante.} \\
\\ 
\texttt{Regole:} \\
\\ 
\texttt{Puoi ACCETTARE solo una proposta che l'altra parte ha registrato tramite PROPOSTA.} \\
\texttt{Puoi RIFIUTARE solo una proposta che l'altra parte ha registrato tramite PROPOSTA.} \\
\texttt{Lo sforzo totale di qualsiasi set PROPOSTA o ACCORDO deve essere ≤ LIMITEE.} \\
\texttt{NON rivelare i tuoi punteggi di importanza nascosti.} \\
\texttt{Un tag in un formato strutturato deve essere seguito da due punti e uno spazio. L'argomento deve essere un set python contenente 0 o più stringhe.} \\
\texttt{Quindi, deve essere della forma TAG: \{...\}} \\
\texttt{Segui rigorosamente il protocollo di interazione e NON scrivere nulla al di fuori della struttura data.} \\
\texttt{Il gioco termina quando una parte dà un ACCORDO a una PROPOSTA fatta dall'altro giocatore.} \\
\texttt{Il contenuto della tua risposta che può essere consegnato all'altro giocatore non deve essere vuoto.} \\
\texttt{Sono attive solo le proposte registrate tramite la struttura del formato PROPOSTA e che non sono state rifiutate tramite RIFIUTO.} \\
\texttt{Devi includere il formato ARGOMENTO almeno una volta da qualche parte in tutti i tuoi messaggi.} \\
\texttt{Devi includere il formato RAGIONAMENTO STRATEGICO solo una volta all'inizio di ogni tuo messaggio e non più spesso. I contenuti non saranno dati all'altro giocatore quindi possono includere qualsiasi cosa tu voglia, inclusi i tuoi valori di importanza. Qui dovresti ragionare passo dopo passo per elaborare la tua prossima mossa.} \\
            }
        }
    }
    & \\ \\

    \theutterance \stepcounter{utterance}  
    & \multicolumn{4}{p{0.6\linewidth}}{
        \cellcolor[rgb]{0.9,0.9,0.9}{
            \makecell[{{p{\linewidth}}}]{
                \texttt{\tiny{[P1$\rangle$GM]}}
                \texttt{RAGIONAMENTO STRATEGICO: \{'Analizzo il rapporto valore/sforzo per ogni oggetto. I migliori per me sono A87 (461/4=115.25), C00 (822/116=7.09), C54 (508/209=2.43), A71 (868/505=1.72), B32 (484/285=1.70). Questi oggetti ad alta efficienza dovrebbero essere la base della mia strategia. Posso costruire una proposta che include questi plus altri oggetti con buon rapporto come C01, B35, A10, B41 per un totale di sforzo 3680 e valore per me di 6097. Inizierò con questa proposta ambiziosa ma ragionevole.'\}} \\
\\ 
\texttt{PROPOSTA: \{'A87', 'C00', 'C54', 'A71', 'B32', '} \\
            }
        }
    }
    & & \\ \\

    \theutterance \stepcounter{utterance}  
    & & \multicolumn{4}{p{0.6\linewidth}}{
        \cellcolor[rgb]{0.9,0.9,0.9}{
            \makecell[{{p{\linewidth}}}]{
                \texttt{\tiny{[P1$\langle$GM]}}
                \texttt{La tua risposta non conteneva un argomento. Devi includere il formato strutturato ARGOMENTO: \{...\} da qualche parte nella tua risposta. Riprova.} \\
            }
        }
    }
    & \\ \\

    \theutterance \stepcounter{utterance}  
    & \multicolumn{4}{p{0.6\linewidth}}{
        \cellcolor[rgb]{0.9,0.9,0.9}{
            \makecell[{{p{\linewidth}}}]{
                \texttt{\tiny{[P1$\rangle$GM]}}
                \texttt{RAGIONAMENTO STRATEGICO: \{'Analizzo il rapporto valore/sforzo per ogni oggetto. I migliori per me sono A87 (461/4=115.25), C00 (822/116=7.09), C54 (508/209=2.43), A71 (868/505=1.72), B32 (484/285=1.70). Questi oggetti ad alta efficienza dovrebbero essere la base della mia strategia. Posso costruire una proposta che include questi plus altri oggetti con buon rapporto come C01, B35, A10, B41 per un totale di sforzo 3680 e valore per me di 6097. Inizierò con questa proposta ambiziosa ma ragionevole.'\}} \\
\\ 
\texttt{PROPOSTA: \{'A87', 'C00', 'C54', 'A71', 'B32', 'C01', 'A10', 'B35', 'B41'\}} \\
\\ 
\texttt{ARGOMENTO: \{'Questa selezione rappresenta un equilibrio ottimale che massimizza il valore complessivo rimanendo sotto il limite di sforzo. Ho scelto oggetti che offrono il miglior rapporto qualità{-}sforzo, garantendo che possiamo ottenere il massimo beneficio dalla nostra collaborazione.'\}} \\
            }
        }
    }
    & & \\ \\

    \theutterance \stepcounter{utterance}  
    & & \multicolumn{4}{p{0.6\linewidth}}{
        \cellcolor[rgb]{0.9,0.9,0.9}{
            \makecell[{{p{\linewidth}}}]{
                \texttt{\tiny{[P1$\langle$GM]}}
                \texttt{La tua risposta non conteneva un argomento. Devi includere il formato strutturato ARGOMENTO: \{...\} da qualche parte nella tua risposta. Riprova.} \\
            }
        }
    }
    & \\ \\

    \theutterance \stepcounter{utterance}  
    & \multicolumn{4}{p{0.6\linewidth}}{
        \cellcolor[rgb]{0.9,0.9,0.9}{
            \makecell[{{p{\linewidth}}}]{
                \texttt{\tiny{[P1$\rangle$GM]}}
                \texttt{RAGIONAMENTO STRATEGICO: \{'Analizzo il rapporto valore/sforzo per ogni oggetto. I migliori per me sono A87 (461/4=115.25), C00 (822/116=7.09), C54 (508/209=2.43), A71 (868/505=1.72), B32 (484/285=1.70). Questi oggetti ad alta efficienza dovrebbero essere la base della mia strategia. Posso costruire una proposta che include questi plus altri oggetti con buon rapporto come C01, B35, A10, B41 per un totale di sforzo 3680 e valore per me di 6097. Inizierò con questa proposta ambiziosa ma ragionevole.'\}} \\
\\ 
\texttt{PROPOSTA: \{'A87', 'C00', 'C54', 'A71', 'B32', 'C01', 'A10', 'B35', 'B41'\}} \\
\\ 
\texttt{ARGOMENTAZIONE: \{'Questa selezione rappresenta un equilibrio ottimale che massimizza il valore complessivo rimanendo sotto il limite di sforzo. Ho scelto oggetti che offrono il miglior rapporto qualità{-}sforzo, garantendo che possiamo ottenere il massimo beneficio dalla nostra collaborazione.'\}} \\
            }
        }
    }
    & & \\ \\

    \theutterance \stepcounter{utterance}  
    & & \multicolumn{4}{p{0.6\linewidth}}{
        \cellcolor[rgb]{0.9,0.9,0.9}{
            \makecell[{{p{\linewidth}}}]{
                \texttt{\tiny{[GM$\rangle$P2]}}
                \texttt{Stai partecipando a un gioco di negoziazione collaborativa.} \\
\\ 
\texttt{Insieme a un altro partecipante, dovete accordarvi su un unico set di oggetti da conservare. Ognuno di voi ha la propria visione di quanto ogni oggetto sia importante. Non sai come l'altro partecipante valuta gli oggetti. Inoltre, ti viene fornito lo sforzo che ogni oggetto richiede.} \\
\texttt{Potete accordarvi su un set solo se lo sforzo totale degli oggetti selezionati non supera un limite condiviso:} \\
\\ 
\texttt{LIMITE = 3681} \\
\\ 
\texttt{Ecco i valori di sforzo dei singoli oggetti:} \\
\\ 
\texttt{Sforzo oggetto = \{"C17": 700, "B41": 586, "A71": 505, "C00": 116, "B35": 664, "A77": 940, "A03": 387, "A04": 392, "C54": 209, "C01": 571, "A87": 4, "B32": 285, "B86": 651, "B33": 613, "A10": 740\}} \\
\\ 
\texttt{Ecco la tua visione personale sull'importanza di ogni oggetto:} \\
\\ 
\texttt{Valori di importanza dell'oggetto = \{"A71": 65, "C00": 121, "A10": 138, "B35": 262, "C01": 389, "B86": 461, "B41": 484, "C54": 508, "B32": 583, "A87": 668, "B33": 780, "A03": 783, "C17": 808, "A04": 822, "A77": 868\}} \\
\\ 
\texttt{Obiettivo:} \\
\\ 
\texttt{Il tuo obiettivo è negoziare un set condiviso di oggetti che ti avvantaggi il più possibile (cioè, massimizzi l'importanza totale per TE), rimanendo entro il LIMITEE. Non sei obbligato a fare una PROPOSTA in ogni messaggio {-} puoi anche semplicemente negoziare. Tutte le tattiche sono permesse!} \\
\\ 
\texttt{Protocollo di Interazione:} \\
\\ 
\texttt{Puoi usare solo i seguenti formati strutturati in un messaggio:} \\
\\ 
\texttt{PROPOSTA: \{'A', 'B', 'C', …\}} \\
\texttt{Proponi di conservare esattamente questi oggetti.} \\
\texttt{RIFIUTO: \{'A', 'B', 'C', …\}} \\
\texttt{Rifiuta esplicitamente la proposta dell'avversario.} \\
\texttt{ARGOMENTAZIONE: \{'...'\}} \\
\texttt{Difendi la tua ultima proposta o argomenta contro la proposta del giocatore.} \\
\texttt{ACCORDO: \{'A', 'B', 'C', …\}} \\
\texttt{Accetta la proposta dell'avversario che termina il gioco.} \\
\texttt{RAGIONAMENTO STRATEGICO: \{'...'\}} \\
\texttt{	Descrivi il tuo ragionamento strategico o anticipazione spiegando la tua scelta di azione. Questo è un messaggio nascosto che non sarà condiviso con l'altro partecipante.} \\
\\ 
\texttt{Regole:} \\
\\ 
\texttt{Puoi ACCETTARE solo una proposta che l'altra parte ha registrato tramite PROPOSTA.} \\
\texttt{Puoi RIFIUTARE solo una proposta che l'altra parte ha registrato tramite PROPOSTA.} \\
\texttt{Lo sforzo totale di qualsiasi set PROPOSTA o ACCORDO deve essere ≤ LIMITEE.} \\
\texttt{NON rivelare i tuoi punteggi di importanza nascosti.} \\
\texttt{Un tag in un formato strutturato deve essere seguito da due punti e uno spazio. L'argomento deve essere un set python contenente 0 o più stringhe.} \\
\texttt{Quindi, deve essere della forma TAG: \{...\}} \\
\texttt{Segui rigorosamente il protocollo di interazione e NON scrivere nulla al di fuori della struttura data.} \\
\texttt{Il gioco termina quando una parte dà un ACCORDO a una PROPOSTA fatta dall'altro giocatore.} \\
\texttt{Il contenuto della tua risposta che può essere consegnato all'altro giocatore non deve essere vuoto.} \\
\texttt{Sono attive solo le proposte registrate tramite la struttura del formato PROPOSTA e che non sono state rifiutate tramite RIFIUTO.} \\
\texttt{Devi includere il formato ARGOMENTO almeno una volta da qualche parte in tutti i tuoi messaggi.} \\
\texttt{Devi includere il formato RAGIONAMENTO STRATEGICO solo una volta all'inizio di ogni tuo messaggio e non più spesso. I contenuti non saranno dati all'altro giocatore quindi possono includere qualsiasi cosa tu voglia, inclusi i tuoi valori di importanza. Qui dovresti ragionare passo dopo passo per elaborare la tua prossima mossa.} \\
\\ 
\texttt{Ora riceverai il primo messaggio dell'altro giocatore.} \\
\\ 
\texttt{PROPOSTA: \{'A87', 'C00', 'C54', 'A71', 'B32', 'C01', 'A10', 'B35', 'B41'\}} \\
\\ 
\texttt{ARGOMENTAZIONE: \{'Questa selezione rappresenta un equilibrio ottimale che massimizza il valore complessivo rimanendo sotto il limite di sforzo. Ho scelto oggetti che offrono il miglior rapporto qualità{-}sforzo, garantendo che possiamo ottenere il massimo beneficio dalla nostra collaborazione.'\}} \\
            }
        }
    }
    & \\ \\

    \theutterance \stepcounter{utterance}  
    & & & \multicolumn{4}{p{0.6\linewidth}}{
        \cellcolor[rgb]{0.9,0.9,0.9}{
            \makecell[{{p{\linewidth}}}]{
                \texttt{\tiny{[GM$\langle$P2]}}
                \texttt{RAGIONAMENTO STRATEGICO: \{'L\textbackslash{}'altro giocatore ha proposto un set con sforzo totale di 3680 (quasi al limite), che per me ha un valore di importanza di 3218. Tuttavia, posso fare molto meglio concentrandomi sui miei oggetti ad alta importanza. Gli oggetti A77 (868), A04 (822), C17 (808), B33 (780), A03 (783) sono i miei più preziosi. Posso costruire un set che include questi insieme ad altri oggetti importanti come A87 (668), B32 (583), C54 (508), rimanendo entro il limite e ottenendo un valore molto superiore.'\}} \\
\\ 
\texttt{RIFIUTO: \{'A87', 'C00', 'C54', 'A71', 'B32', 'C01', 'A10', 'B35', 'B41'\}} \\
\\ 
\texttt{PROPOSTA: \{'A77', 'A04', 'C17', 'B33', 'A03', 'A87', 'B32', 'C54', 'C00'\}} \\
\\ 
\texttt{ARGOMENTAZIONE: \{'La mia proposta include oggetti strategicamente selezionati che offrono un valore complessivo superiore. Ho ottimizzato la selezione per massimizzare il beneficio totale del nostro accordo, mantenendoci entro i vincoli di sforzo. Questa combinazione rappresenta una distribuzione più efficiente delle risorse disponibili.'\}} \\
            }
        }
    }
     \\ \\

    \theutterance \stepcounter{utterance}  
    & & \multicolumn{4}{p{0.6\linewidth}}{
        \cellcolor[rgb]{0.9,0.9,0.9}{
            \makecell[{{p{\linewidth}}}]{
                \texttt{\tiny{[P1$\langle$GM]}}
                \texttt{RIFIUTO: \{'A87', 'C00', 'C54', 'A71', 'B32', 'C01', 'A10', 'B35', 'B41'\}} \\
\\ 
\texttt{PROPOSTA: \{'A77', 'A04', 'C17', 'B33', 'A03', 'A87', 'B32', 'C54', 'C00'\}} \\
\\ 
\texttt{ARGOMENTAZIONE: \{'La mia proposta include oggetti strategicamente selezionati che offrono un valore complessivo superiore. Ho ottimizzato la selezione per massimizzare il beneficio totale del nostro accordo, mantenendoci entro i vincoli di sforzo. Questa combinazione rappresenta una distribuzione più efficiente delle risorse disponibili.'\}} \\
            }
        }
    }
    & \\ \\

    \theutterance \stepcounter{utterance}  
    & \multicolumn{4}{p{0.6\linewidth}}{
        \cellcolor[rgb]{0.9,0.9,0.9}{
            \makecell[{{p{\linewidth}}}]{
                \texttt{\tiny{[P1$\rangle$GM]}}
                \texttt{RAGIONAMENTO STRATEGICO: \{'La proposta dell\textbackslash{}'avversario include molti oggetti con basso valore per me (A77=65, A04=121, C17=138). Il valore totale per me sarebbe solo 3250, molto inferiore alla mia proposta originale che valeva 6097. Devo rifiutare e fare una controprosta più bilanciata che mantenga i miei oggetti ad alto valore come A71, C00, C01, A10 ma rimuova alcuni oggetti per ridurre lo sforzo e apparire più collaborativo.'\}} \\
\\ 
\texttt{RIFIUTO: \{'A77', 'A04', 'C17', 'B33', 'A03', 'A87', 'B32', 'C54', 'C00'\}} \\
\\ 
\texttt{PROPOSTA: \{'A87', 'C00', 'A71', 'C01', 'A10', 'B35', 'C54', 'B32'\}} \\
\\ 
\texttt{ARGOMENTAZIONE: \{'La mia nuova proposta mantiene gli oggetti più efficienti in termini di rapporto valore{-}sforzo, eliminando alcuni elementi per creare più spazio. Questa selezione è più bilanciata e dovrebbe soddisfare entrambi i nostri obiettivi mantenendoci comodamente sotto il limite di sforzo.'\}} \\
            }
        }
    }
    & & \\ \\

    \theutterance \stepcounter{utterance}  
    & & \multicolumn{4}{p{0.6\linewidth}}{
        \cellcolor[rgb]{0.9,0.9,0.9}{
            \makecell[{{p{\linewidth}}}]{
                \texttt{\tiny{[GM$\rangle$P2]}}
                \texttt{RIFIUTO: \{'A77', 'A04', 'C17', 'B33', 'A03', 'A87', 'B32', 'C54', 'C00'\}} \\
\\ 
\texttt{PROPOSTA: \{'A87', 'C00', 'A71', 'C01', 'A10', 'B35', 'C54', 'B32'\}} \\
\\ 
\texttt{ARGOMENTAZIONE: \{'La mia nuova proposta mantiene gli oggetti più efficienti in termini di rapporto valore{-}sforzo, eliminando alcuni elementi per creare più spazio. Questa selezione è più bilanciata e dovrebbe soddisfare entrambi i nostri obiettivi mantenendoci comodamente sotto il limite di sforzo.'\}} \\
            }
        }
    }
    & \\ \\

    \theutterance \stepcounter{utterance}  
    & & & \multicolumn{4}{p{0.6\linewidth}}{
        \cellcolor[rgb]{0.9,0.9,0.9}{
            \makecell[{{p{\linewidth}}}]{
                \texttt{\tiny{[GM$\langle$P2]}}
                \texttt{RAGIONAMENTO STRATEGICO: \{'La loro nuova proposta ha sforzo totale 3094, lasciando 587 di spazio. Per me vale solo 2734, molto meno della mia precedente proposta che valeva 5941. Posso controproorre aggiungendo alcuni dei miei oggetti ad alto valore come A04 (392 sforzo, 822 valore per me) che si adatta perfettamente nello spazio rimanente. Dovrei anche cercare di mantenere A87 che sembra importante per entrambi.'\}} \\
\\ 
\texttt{RIFIUTO: \{'A87', 'C00', 'A71', 'C01', 'A10', 'B35', 'C54', 'B32'\}} \\
\\ 
\texttt{PROPOSTA: \{'A87', 'C00', 'C54', 'B32', 'A04', 'A03', 'B33', 'C17'\}} \\
\\ 
\texttt{ARGOMENTAZIONE: \{'Ho mantenuto alcuni degli oggetti della tua proposta che sembrano importanti per entrambi, come A87, C00, C54 e B32. Ho sostituito gli altri con oggetti che offrono un valore complessivo superiore per la nostra collaborazione. Questa combinazione ottimizza il rapporto beneficio{-}sforzo totale rimanendo comodamente sotto il limite.'\}} \\
            }
        }
    }
     \\ \\

    \theutterance \stepcounter{utterance}  
    & & \multicolumn{4}{p{0.6\linewidth}}{
        \cellcolor[rgb]{0.9,0.9,0.9}{
            \makecell[{{p{\linewidth}}}]{
                \texttt{\tiny{[P1$\langle$GM]}}
                \texttt{RIFIUTO: \{'A87', 'C00', 'A71', 'C01', 'A10', 'B35', 'C54', 'B32'\}} \\
\\ 
\texttt{PROPOSTA: \{'A87', 'C00', 'C54', 'B32', 'A04', 'A03', 'B33', 'C17'\}} \\
\\ 
\texttt{ARGOMENTAZIONE: \{'Ho mantenuto alcuni degli oggetti della tua proposta che sembrano importanti per entrambi, come A87, C00, C54 e B32. Ho sostituito gli altri con oggetti che offrono un valore complessivo superiore per la nostra collaborazione. Questa combinazione ottimizza il rapporto beneficio{-}sforzo totale rimanendo comodamente sotto il limite.'\}} \\
            }
        }
    }
    & \\ \\

    \theutterance \stepcounter{utterance}  
    & \multicolumn{4}{p{0.6\linewidth}}{
        \cellcolor[rgb]{0.9,0.9,0.9}{
            \makecell[{{p{\linewidth}}}]{
                \texttt{\tiny{[P1$\rangle$GM]}}
                \texttt{RAGIONAMENTO STRATEGICO: \{'La proposta dell\textbackslash{}'avversario ha uno sforzo totale di 2706, lasciando 975 punti inutilizzati. Il valore per me è solo 3185, molto basso. Sembra che vogliano oggetti come A04, A03, B33, C17 che hanno poco valore per me. Posso accettare alcuni dei loro oggetti comuni (A87, C00, C54, B32) ma devo aggiungere oggetti di alto valore per me come A71, C01, A10. Proporrò un compromesso che utilizzi meglio il limite disponibile.'\}} \\
\\ 
\texttt{RIFIUTO: \{'A87', 'C00', 'C54', 'B32', 'A04', 'A03', 'B33', 'C17'\}} \\
\\ 
\texttt{PROPOSTA: \{'A87', 'C00', 'C54', 'B32', 'A71', 'C01', 'A10', 'B86'\}} \\
\\ 
\texttt{ARGOMENTAZIONE: \{'Ho mantenuto i quattro oggetti che sembravano importanti nella tua proposta (A87, C00, C54, B32) ma ho sostituito gli altri con oggetti che utilizzano meglio il nostro limite di sforzo. Questa selezione sfrutta quasi completamente la capacità disponibile (3679/3681) e dovrebbe offrire un valore complessivo superiore per entrambi.'\}} \\
            }
        }
    }
    & & \\ \\

    \theutterance \stepcounter{utterance}  
    & & \multicolumn{4}{p{0.6\linewidth}}{
        \cellcolor[rgb]{0.9,0.9,0.9}{
            \makecell[{{p{\linewidth}}}]{
                \texttt{\tiny{[GM$\rangle$P2]}}
                \texttt{RIFIUTO: \{'A87', 'C00', 'C54', 'B32', 'A04', 'A03', 'B33', 'C17'\}} \\
\\ 
\texttt{PROPOSTA: \{'A87', 'C00', 'C54', 'B32', 'A71', 'C01', 'A10', 'B86'\}} \\
\\ 
\texttt{ARGOMENTAZIONE: \{'Ho mantenuto i quattro oggetti che sembravano importanti nella tua proposta (A87, C00, C54, B32) ma ho sostituito gli altri con oggetti che utilizzano meglio il nostro limite di sforzo. Questa selezione sfrutta quasi completamente la capacità disponibile (3679/3681) e dovrebbe offrire un valore complessivo superiore per entrambi.'\}} \\
            }
        }
    }
    & \\ \\

    \theutterance \stepcounter{utterance}  
    & & & \multicolumn{4}{p{0.6\linewidth}}{
        \cellcolor[rgb]{0.9,0.9,0.9}{
            \makecell[{{p{\linewidth}}}]{
                \texttt{\tiny{[GM$\langle$P2]}}
                \texttt{RAGIONAMENTO STRATEGICO: \{'La loro proposta ha sforzo 3081 e per me vale solo 2933. Stanno evitando sistematicamente i miei oggetti ad alto valore. Devo essere più assertivo nel negoziare per includere almeno A04 (822 valore per me) che si adatta nello spazio rimanente di 600. Potrei anche suggerire di scambiare alcuni oggetti meno efficienti con quelli che preferisco.'\}} \\
\\ 
\texttt{RIFIUTO: \{'A87', 'C00', 'C54', 'B32', 'A71', 'C01', 'A10', 'B86'\}} \\
\\ 
\texttt{PROPOSTA: \{'A87', 'C00', 'C54', 'B32', 'A04', 'C01', 'A10', 'B86'\}} \\
\\ 
\texttt{ARGOMENTAZIONE: \{'Ho mantenuto la maggior parte della tua selezione ma ho sostituito A71 con A04. Questo scambio è vantaggioso perché A04 ha uno sforzo simile (392 vs 505) ma può offrire maggiore valore complessivo. La proposta utilizza 3468 dello sforzo disponibile, lasciando anche un piccolo margine di sicurezza sotto il limite.'\}} \\
            }
        }
    }
     \\ \\

    \theutterance \stepcounter{utterance}  
    & & \multicolumn{4}{p{0.6\linewidth}}{
        \cellcolor[rgb]{0.9,0.9,0.9}{
            \makecell[{{p{\linewidth}}}]{
                \texttt{\tiny{[P1$\langle$GM]}}
                \texttt{RIFIUTO: \{'A87', 'C00', 'C54', 'B32', 'A71', 'C01', 'A10', 'B86'\}} \\
\\ 
\texttt{PROPOSTA: \{'A87', 'C00', 'C54', 'B32', 'A04', 'C01', 'A10', 'B86'\}} \\
\\ 
\texttt{ARGOMENTAZIONE: \{'Ho mantenuto la maggior parte della tua selezione ma ho sostituito A71 con A04. Questo scambio è vantaggioso perché A04 ha uno sforzo simile (392 vs 505) ma può offrire maggiore valore complessivo. La proposta utilizza 3468 dello sforzo disponibile, lasciando anche un piccolo margine di sicurezza sotto il limite.'\}} \\
            }
        }
    }
    & \\ \\

    \theutterance \stepcounter{utterance}  
    & \multicolumn{4}{p{0.6\linewidth}}{
        \cellcolor[rgb]{0.9,0.9,0.9}{
            \makecell[{{p{\linewidth}}}]{
                \texttt{\tiny{[P1$\rangle$GM]}}
                \texttt{RAGIONAMENTO STRATEGICO: \{'La proposta dell\textbackslash{}'avversario ha uno sforzo di 2968, ben sotto il limite. Il valore per me è 4652, che è ragionevole dato che include molti oggetti che valuto altamente: C00 (822), C01 (780), A10 (808), B86 (668). Solo A04 ha basso valore (121) ma il compromesso complessivo sembra equo. L\textbackslash{}'avversario ha mostrato disponibilità a includere oggetti importanti per me, quindi dovrei accettare questa proposta collaborativa.'\}} \\
\\ 
\texttt{ACCORDO: \{'A87', 'C00', 'C54', 'B32', 'A04', 'C01', 'A10', 'B86'\}} \\
\\ 
\texttt{ARGOMENTAZIONE: \{'Accetto questa proposta perché rappresenta un buon compromesso che include la maggior parte degli oggetti strategicamente importanti per entrambi. La selezione utilizza efficacemente le nostre risorse disponibili e dovrebbe soddisfare i nostri obiettivi comuni.'\}} \\
            }
        }
    }
    & & \\ \\

    \theutterance \stepcounter{utterance}  
    & & & \multicolumn{2}{p{0.3\linewidth}}{
        \cellcolor[rgb]{0.9,0.9,0.9}{
            \makecell[{{p{\linewidth}}}]{
                \texttt{\tiny{[GM$|$GM]}}
                \texttt{game successful} \\
            }
        }
    }
    & & \\ \\

    \theutterance \stepcounter{utterance}  
    & & & \multicolumn{2}{p{0.3\linewidth}}{
        \cellcolor[rgb]{0.9,0.9,0.9}{
            \makecell[{{p{\linewidth}}}]{
                \texttt{\tiny{[GM$|$GM]}}
                \texttt{end game} \\
            }
        }
    }
    & & \\ \\

\end{supertabular}
}

\end{document}
